\section{Overview}
\label{sec:overview}

KeyQuorum is a secure file-sharing system centered on hardware key sharing.
Files are encrypted and bound to registered physical tokens, such as USB
devices. Access requires the necessary hardware keys to be presented before
a protected file can be unlocked, providing layered, hardware-backed access
control.

\begin{noteBox}
This physical binding depends on a key actually being placed in a
\code{keyquorum-device} container (Section~\ref{sec:devices-and-custody}).
An unplaced \flag{--share-file} key file is classified by its own key
identity: two such files can both count toward
\code{minimum_physical_devices} even while sitting on the same physical
drive, since nothing ties either one to a container's \code{device_id}.
\end{noteBox}

\subsection{Concept}

\begin{itemize}
  \item A \textbf{key} --- the data key protecting a file, or a secret split
  for its own sake --- can be divided recursively: split into parts, any of
  which can itself be split further, forming a tree. A company master key
  can split across departments, and one department's own share can split
  again across its team. A flat ``$M$-of-$N$ hardware keys'' quorum is just
  the simplest one-level tree.
  \item \textbf{Unlocking} a protected file means reconstructing its key's
  tree: presenting a quorum of the registered keys, rather than a single
  token or password.
  \item The goal is layered, hardware-backed access control that resists
  single-point compromise. A lost or stolen key alone should not be enough
  to unlock protected data.
\end{itemize}

\subsection{What this manual covers}

This manual documents the \code{keyquorum} and \code{keyquorum-device}
command-line tools: generating and registering hardware keys, building and
reconstructing quorum trees, protecting files and credentials, private
signing bridges, authenticated updates to a live tree, moving identities
between physical devices, sending files directly to another person, tracked
files that carry their own signed revision history, and the mailbox relay that carries all of the above between machines that are
not open at the same time.

\begin{noteBox}
Passwords, PINs, and hardware-key material are always prompted for
interactively, or read from a \flag{--share-file} key path, rather than
taken as plain command-line arguments. This keeps secrets out of shell
history and process listings.
\end{noteBox}

\subsection{Status}

The CLI covers most of what this manual describes. Hardware-key quorum
splitting and reconstruction are implemented in software: a key file with no
container placement is still one device, several identities can share one
container, a non-root tree restructure waits for its parent to countersign,
and \code{transfer copy} / \code{transfer move} carry an active identity to
a second device that is open at the same time, or through the mailbox as a
sealed letter when it is not.

Tracked files (section~\ref{sec:file-history}) are implemented through
\code{keyquorum file}; that section ends with what their design still
plans. Two other things remain outside the current scope: persisted private-key custody
beyond a key file or a \code{keyquorum-device} container (for example an OS
keychain), and \code{import} of password-vault or locked-file
\code{export} bundles (the bundle format itself is final; private-bridge
\code{.kqpb} import is already implemented).

\subsection{Browser lab}

KeyQuorum Lab runs this crate in a browser as a small sandboxed machine,
compiled to WebAssembly. It has its own files, a SQLite store per person,
mock USB drives, and a relay answering in the page. Every button and every
terminal line runs the real \code{keyquorum} and \code{keyquorum-device}
commands, so the lab demonstrates the same behavior described in this
manual. It is itself a browser-based virtual machine that emulates
USB-token custody entirely in memory, with published default passphrases
--- so it is not physical hardware-backed security. It is a good place to try the
command sequences in this manual without provisioning real hardware.

\begin{noteBox}
The lab's own GUI --- its tabs, the Active user bar, and its guided
in-browser tutorials --- is documented separately in the \emph{KeyQuorum
Lab} manual, not in this one. This manual covers the \code{keyquorum} /
\code{keyquorum-device} CLI itself, which the lab runs underneath its
interface unchanged.
\end{noteBox}
